% ============================================================
%  Chương VII – Thách thức và xu hướng
% ============================================================
\chapter{Thách thức và xu hướng}
\label{chap:challenges}

PKI vẫn là nền tảng để ràng buộc public key với danh tính.
Khi số lượng certificate tăng, rủi ro chính chuyển sang kiểm soát
vòng đời, niềm tin và khóa riêng.

% ----------------------------------------------------------
\section{Thách thức hiện tại}
\label{sec:7.1}

Thách thức đầu tiên là quản lý vòng đời chứng chỉ. Một certificate có thể
được cấp đúng nhưng gây lỗi khi hết hạn, không được gia hạn, hoặc bị bỏ quên
trong một trust store cũ. Ở quy mô lớn, tổ chức phải kiểm kê certificate, chủ
sở hữu, CA phát hành, thời hạn và mục đích sử dụng. Thiếu kiểm kê làm xoay
vòng khóa và thay thế certificate trở thành thao tác thủ công dễ sai.

Thu hồi chứng chỉ cũng có giới hạn thực tế. CRL được định nghĩa trong profile
X.509, nhưng có thể lớn và có độ trễ phân phối \cite{RFC5280}. OCSP cho phép
kiểm tra trạng thái trực tuyến, nhưng lại phụ thuộc vào responder và chính
sách xử lý khi dịch vụ trạng thái không sẵn sàng \cite{RFC6960}. Vì vậy, thu
hồi không phải cơ chế tức thời tuyệt đối.

Niềm tin vào CA là rủi ro tập trung. Nếu CA cấp sai hoặc bị xâm phạm, trình
duyệt và hệ điều hành phải dựa vào root store, chính sách CA và cơ chế distrust
để giới hạn thiệt hại. CA/B Forum Baseline Requirements đặt yêu cầu cho CA
publicly-trusted \cite{CABForum}.

% ----------------------------------------------------------
\section{Tự động hóa và giám sát}
\label{sec:7.2}

Xu hướng chính là giảm thao tác thủ công. ACME chuẩn hóa cách client tạo
order, chứng minh quyền kiểm soát identifier, gửi CSR và nhận certificate
\cite{RFC8555}. Certificate lifecycle management mở rộng ý tưởng này sang
kiểm kê, cấp phát, gia hạn, thu hồi và audit. Short-lived certificate cũng
giảm khoảng thời gian rủi ro khi khóa hoặc certificate bị lộ.

Tự động hóa cần đi kèm policy. Nếu hệ thống cấp certificate không kiểm soát
được ai có quyền yêu cầu, tên nào được cấp, hoặc profile nào được dùng, lỗi
thủ công chỉ chuyển thành lỗi tự động. Vì vậy, inventory, phân quyền và audit
là điều kiện vận hành, không phải phần phụ.

Bảng~\ref{tab:pki-trends} tóm tắt vai trò của một số hướng xử lý hiện nay.

\begin{table}
    \centering
    \small
    \caption{Một số hướng xử lý thách thức PKI}
    \label{tab:pki-trends}
    \begin{tabularx}{\textwidth}{@{} >{\raggedright\arraybackslash}p{0.22\textwidth}
        >{\raggedright\arraybackslash}p{0.30\textwidth}
        >{\raggedright\arraybackslash}X @{} }
        \toprule
        \textbf{Hướng} & \textbf{Giảm rủi ro} & \textbf{Giới hạn} \\
        \midrule
        ACME và CLM
            & Giảm lỗi cấp phát và gia hạn thủ công
            & Cần policy, phân quyền và inventory. \\
        Short-lived certificate
            & Giảm thời gian hiệu lực khi có sự cố
            & Phụ thuộc tự động hóa gia hạn ổn định. \\
        Certificate Transparency
            & Tăng khả năng phát hiện certificate cấp sai
            & Không thay thế validation hoặc revocation. \\
        Crypto-agility
            & Cho phép thay thuật toán và profile
            & Cần hỗ trợ từ CA, HSM, TLS stack và client. \\
        \bottomrule
    \end{tabularx}
\end{table}

Certificate Transparency bổ sung lớp giám sát cho Web PKI. RFC~9162 mô tả CT
bằng các log append-only dựa trên Merkle Tree, giúp certificate đã cấp có thể
được audit công khai \cite{RFC9162}. CT làm tăng minh bạch và hỗ trợ phát
hiện mis-issuance. Nó không tự xác thực domain, không thu hồi certificate và
không thay thế chính sách CA.

% ----------------------------------------------------------
\section{PKI trong hệ thống hiện đại}
\label{sec:7.3}

PKI ngày càng gắn với machine identity. Service, container, workload và thiết
bị IoT đều cần danh tính mật mã để kết nối an toàn. Khác với người dùng,
những thực thể này có số lượng lớn, vòng đời ngắn và thường được tạo tự động.
Điều đó làm certificate inventory và policy trở thành phần của hạ tầng vận
hành.

Trong kiến trúc zero trust và microservice, mTLS dùng certificate để xác thực
hai chiều giữa các service. Cách làm này giảm phụ thuộc vào vị trí mạng,
nhưng tăng yêu cầu đối với CA nội bộ, phân phối trust anchor và xoay vòng
private key. Internal PKI vì vậy cần tích hợp với quản trị danh tính, secret
management và audit thay vì chỉ phát hành certificate rời rạc.

% ----------------------------------------------------------
\section{Post-quantum}
\label{sec:7.4}

Máy lượng tử đủ mạnh sẽ đe dọa RSA, Diffie-Hellman, ECDSA và ECDH vì các hệ
này dựa trên bài toán phân tích số hoặc logarit rời rạc. NIST đã công bố các
chuẩn hậu lượng tử chính năm 2024, gồm KEM và thuật toán chữ ký số
\cite{NIST-PQC2024}. Tác động tới PKI không chỉ là thay thuật toán. Certificate
và chữ ký có thể lớn hơn, chain validation phức tạp hơn, còn CA, HSM, TLS
stack và trust store phải cùng hỗ trợ profile mới.

Vì chưa nên dự đoán chắc mốc thời gian cho khả năng phá vỡ lượng tử, hướng
thực tế là crypto-agility. Một hệ PKI cần thay được thuật toán, profile và
policy mà không làm gián đoạn cấp phát, xác minh và thu hồi certificate. Đây
là lý do PKI phải tiến hóa ở tầng vận hành, không chỉ ở lựa chọn thuật toán.
